\documentclass[10pt]{article}
\usepackage[margin=0.9in]{geometry}
\usepackage[T1]{fontenc}
\usepackage{mathptmx}
\usepackage{microtype}
\usepackage{amsmath}
\usepackage{xcolor}
\usepackage{enumitem}
\usepackage[colorlinks=true,urlcolor=blue!50!black,linkcolor=blue!50!black]{hyperref}
\newenvironment{rev}{\par\medskip\noindent\color{black!70}\itshape}{\par}
\newcommand{\resp}{\par\smallskip\noindent\textbf{Response.} }
\newcommand{\chg}{\par\smallskip\noindent\textbf{Changes.} }
\setlist{nosep}
\title{\textbf{Response to reviewers}\\\large Round 2}
\author{Kyler Zook}
\date{October 2026}
\begin{document}
\maketitle

We thank the reviewer again. Every point was correct, and each is addressed below. Two new analyses changed what we can claim:
\begin{itemize}
\item the positive control's sensitivity depends strongly on the generator, so the ``exclusion'' language is withdrawn;
\item calibration shows that the earlier negative magnitude scores were largely amplitude miscalibration.
\end{itemize}
Revision-history phrasing has been removed from the manuscript; it is kept here and in the round-1 letter (\texttt{paper/response\_round1.pdf}).

\section*{1. Positive control and the ``$\ge5\%$ excluded'' claim}
\resp
\begin{itemize}
\item \textbf{Definition of $\alpha$.} The residual component is now orthogonalised against the cascade direction, so $\alpha$ is exactly the fraction of the noise-free target's signal energy along the cascade direction. The definition is given in Section~2.3, and realised fractions are reported for every run. The original non-orthogonal construction had realised fractions 0.005, 0.053, 0.102, 0.250 and 0.499 for nominal 0, 0.05, 0.1, 0.25 and 0.5.
\item \textbf{Robustness.} We evaluated mediator ranking errors (true mediators drawn at random), mediators whose coordinates are not visible in the target, heterogeneous mediator strengths, halved reliability of either screen, and a combination (Supplementary Table~S13, Fig.~3c).
\begin{itemize}
\item Heterogeneity and unequal reliability reduced sensitivity but did not remove it.
\item Random mediator ranking reduced it sharply.
\item Invisible mediators, and the combined adverse generator, removed it entirely, even at $\alpha=0.5$.
\end{itemize}
\item \textbf{Wording.} We did not attempt a calibrated exclusion procedure, because these results show that power depends on untestable generator properties. We adopted the reviewer's wording: the pipeline detects a nominal 5\% injected component under the tested semi-synthetic generator, including reduced mediator coverage, provided the mediators are among the knockdown's strongest measured transcript responders. Cascades through weakly responding or transcriptionally invisible mediators are explicitly not excluded.
\end{itemize}
\chg Abstract, Section~2.3, Table~5 (base generator, exact $\alpha$), Fig.~3c, Supplementary Table~S13, Discussion, Limitations.

\section*{2. Context categories}
\resp Correct; the earlier text omitted 23.6\% (Norman) and 36.3\% (Wessels) of significant changes. We added an ``other (intermediate strength)'' category and defined all five categories explicitly. Each significant change is assigned to exactly one, and the five sum to the total (Table~7, Fig.~5).
Reversal percentages now state their denominator: 1.8\% and 1.1\% of all significant changes; 1.4\% (218 of 15{,}778) and 1.1\% (18 of 1{,}657) of strong responses shared by both contexts.
We state that $|z_\text{diff}|>3$ is a nominal per-gene threshold without multiple-testing control. It is used to compare aggregate counts with the matched null, not to certify individual genes.
\chg Abstract, Section~2.5, Table~7, Fig.~5.

\section*{3. Magnitude scores and calibration}
\resp We added training-only scalar calibration, $\hat y_\text{cal}=\lambda\hat y$ with $\lambda\in[0,10]$ ($\lambda=0$ allowed). $\lambda$ is chosen per fold, predictor and target on leave-one-out predictions for 300 training knockdowns, maximising the mean per-knockdown $R^2_0$ (the reported metric).
\begin{itemize}
\item After calibration all predictors are positive on independent raw profiles (0.033--0.047) and small but positive on mean-centred (0.005--0.013) and specific targets (0.002--0.004). Fitted $\lambda$ values were 0.10--0.46, so the earlier negative values were largely over-scaled predictions.
\item Paired per-knockdown comparisons show that more neighbours helps beyond shrinkage on the specific target: $k=100$ minus $k=10$ gives $\Delta$corr $+0.023$ [0.016, 0.030] and calibrated $\Delta R^2_0$ $+0.0024$ [0.0016, 0.0031].
\item The abstract no longer says that no predictor beats zero; it reports the calibrated value. The round-1 letter's statement was inconsistent with Table~3 (mean profile $R^2_0=0.007$ on independent raw profiles), as the reviewer noted.
\end{itemize}
\chg Abstract, Section~2.2, Supplementary Tables~S14--S15.

\section*{4. Noise check}
\resp The 1.30 and 1.18 ratios are now described as consistency checks, not an attribution. The text states that they are medians of ratios over different pair sets, and that screen effects, extra perturbed-profile noise and true between-screen differences cannot be separated.
We added the assumption behind allocating noise by dimensionality (isotropic noise), and note that axes fitted to the same noisy matrix absorb extra noise, which biases the specific share slightly low.
The headline noise fractions now carry estimator sensitivity (pooled 0.47--0.50, median 0.74--0.78 across five estimators) and a bootstrap over controls (Supplementary Table~S11). The text states that these intervals cover the noise constant only, not noise-model misspecification.
The Introduction now says that cross-screen comparison measures agreement without a noise model, while interpreting disagreement as noise requires assumptions.
\chg Introduction, Section~2.1, Supplementary Table~S11.

\section*{5. Subspace vs individual axes}
\resp The text now says that the fitted three-dimensional subspaces share two closely aligned directions. We also report matched individual-axis cosines: GW axis 1 matches essential axis 1 ($|\cos|=0.95$), axis 2 matches axis 2 (0.91), and axis 3 matches nothing ($\le0.12$). Top-loading genes of the two aligned (canonical) directions are listed in the text.
\chg Section~2.1, Supplementary Table~S12.

\section*{6. Propagation control}
\resp Both operators are now scored against the same targets on the same log-normalised scale: (i) the noise-free effect mapped to log scale, and (ii) the measured effect. True-operator predictions use the same transform.
For regulators with depth $\ge3$, the mean of $D=2$--4 minus $D=1$ was:
\begin{itemize}
\item true operator: $+0.19$ [0.16, 0.23] (noise-free target) and $+0.07$ [0.05, 0.09] (measured target);
\item stabilised estimated operator: $+0.04$ [0.01, 0.07] and $-0.02$ [$-0.05$, $-0.00$].
\end{itemize}
We adopted the reviewer's wording: deeper propagation helps with the true operator in this simulator, and the estimated pipeline does not reproduce that benefit. The text also says ``no consistent improvement beyond one step'' rather than ``one step was best''. Residual differences (linearisation, cell-to-cell heterogeneity) are stated.
\chg Supplement S8, Supplementary Table~S7, Fig.~S4.

\section*{7. Reproducibility}
\resp
\begin{itemize}
\item \textbf{Repository.} Code, results, manuscript sources and environment files are at \url{https://github.com/kylerzook/cross-screen-perturbseq}, tagged release \texttt{v1.2}. The tag \texttt{v1.1-revision1} cited in the round-1 letter existed only in a local working copy and is not used. An archived release with a DOI will be added on acceptance.
\item \textbf{Environments.} The main environment (main text and Supplement S5--S10) is pinned in \texttt{requirements.txt}. The environment used for most S1--S4 runs, including PyTorch 2.14.1+cpu, is pinned in \texttt{requirements-s1-s4.txt}, exported from the run machine.
\item \textbf{Clean-setup check.} We created a fresh virtual environment from \texttt{requirements.txt}, cloned the repository, and re-ran the main-text pipeline from the raw files: decomposition, noise checks, binding, Norman preprocessing and context, benchmark, and cascade test. Outputs were compared with the committed results. All 1{,}069 re-computed numeric values in the decomposition, sensitivity, noise, binding, benchmark and cascade outputs matched exactly. The context analysis matched exactly in observed counts, but its null replicates differed within Monte Carlo error: comparison order came from Python set iteration, which depends on the hash seed. We made the order deterministic, re-ran the analysis (Wessels null 562 instead of 565 significant changes; ratio unchanged at 3.0), and verified identical output under different hash seeds (\texttt{docs/13\_clean\_run\_check.md}). The binding step used the cached per-TF promoter features (\texttt{data/processed/test2}), whose generation (S1) was not re-run.
\item \textbf{Scope.} The remaining commands in \texttt{run\_all.sh} are described as supplied, not verified.
\end{itemize}

\section*{Smaller corrections}
\begin{itemize}
\item \textbf{Cell-count tertiles:} the text now distinguishes absolute expected noise ($cG/n$, lowest with many cells) from fractional noise (highest there, because those knockdowns have weaker responses).
\item \textbf{Hub score:} now ``not fully explained by these coarse stratifications''. Moderate correlations with accessibility (0.33) and expression (0.23) are acknowledged.
\item \textbf{Measured ChIP:} restricted to ``no detectable discrimination under this responder definition''.
\item \textbf{Revision history:} phrases (``reviewer-requested'', ``survive the revision'', ``our first version'', ``earlier version'') were removed from the manuscript and Supplement.
\end{itemize}
\end{document}
